% =============================================================================
%  Resumen / Abstract — se redacta AL FINAL, una vez cerrado el resto.
% =============================================================================

\chapter*{Resumen}
\addcontentsline{toc}{chapter}{Resumen}

\todo{Resumen en español, 250–500 palabras. Cubrir:
\begin{itemize}
  \item Problema (IA generativa de imágenes mediante modelos de difusión).
  \item Sistema desarrollado: librería \texttt{diffusion\_lib} con procesos VE/VP, samplers EM/PC/PF-ODE, schedules lineal/coseno, métricas BPD/FID/IS y generación controlable (CFG e imputación).
  \item Resultados principales: FID/IS sobre los datasets entrenados.
  \item Conclusiones de alto nivel.
\end{itemize}}

\vspace{1cm}
\textbf{Palabras clave:} modelos de difusión, IA generativa, score matching, SDE, U-Net, FID, BPD.

\chapter*{Abstract}
\addcontentsline{toc}{chapter}{Abstract}

\todo{Resumen en inglés, 250–500 palabras. Traducción del español.}

\vspace{1cm}
\textbf{Keywords:} diffusion models, generative AI, score matching, SDE, U-Net, FID, BPD.
